\section{Four-edge cuts: exact states and a forest estimate}
\label{sec:four-cut-states}

\subsection{Exact gluing formula}
Let $X,Y$ be connected four-poles with ranks $h,k$.
Their four distinct degree-two ports have matching labels $1,2,3,4$.
All other vertices have degree three.
Join equally labelled ports by four new edges to obtain $G$.
Then $r=\beta(G)=h+k+3$.

Write $p_{ij}=P_{ij}(X)$ and $q_{ij}=P_{ij}(Y)$.
Let $f_\alpha$ count the vertex-disjoint path pairs in $X$ with pairing
$\alpha\in\{12|34,13|24,14|23\}$.
Define $g_\alpha$ similarly in $Y$, and put $F=\sum_\alpha f_\alpha$,
$J=\sum_\alpha g_\alpha$.
The exact cycle decomposition is
\begin{equation}\label{eq:four-cut-exact}
 c(G)=c(X)+c(Y)+\sum_{i<j}p_{ij}q_{ij}
       +FJ-\sum_\alpha f_\alpha g_\alpha.
\end{equation}
Indeed, a crossing cycle uses either two or four cut edges.
Two cut edges give one corresponding terminal path in each shore.
Four cut edges give a path pair in each shore.
Different pairings combine into one cycle, while equal pairings give two cycles.
The half-density target for the crossing contribution is $4\,2^{h+k}$.

\subsection{At most two cycles in a graph with many terminal paths}

\begin{lemma}[Two-component cycle packings]
\label{lem:high-path-two-cycle-packings}
For each $\varepsilon>0$, there is a constant $B_\varepsilon$ with the
following property. Let $N$ be a normalized two-pole of rank $b$ with
\[
 P(N)\ge(1/2+\varepsilon)2^b.
\]
The number of nonempty unions of at most two vertex-disjoint cycles in $N$
is at most
\[
 B_\varepsilon(b+1)^4 2^{b/2}.
\]
\end{lemma}

\begin{proof}
First assume $N$ is bridgeless.
Use the shortest terminal path, distance layers, and initial marked set from
Lemma~\ref{lem:thin-layer-cycle-count}.
The additive stability theorem bounds both $\Delta$ and $H$ by
$C_\varepsilon$. The initial marked set has at most
$6\Delta+4H+2$ vertices.

Choose the minimum and maximum distance layers of each selected cycle.
There are at most $(D+1)^4$ choices for two cycles, where
$D\le b+2\Delta$ is the largest distance.
For one cycle, the corresponding bound is $(D+1)^2$.
Mark the chosen extreme layers when they were previously ordinary.
At most eight vertices are added, so the marked set $B$ satisfies
\[
 |B|\le6\Delta+4H+10.
\]
Fix all selected edges incident with $B$, at a cost of at most $2^{3|B|}$.

Every remaining ordinary layer used by one cycle lies strictly between
that cycle's extreme layers. The cycle therefore uses both vertices, with
two strands through the layer.
Two vertex-disjoint cycles cannot both use an ordinary layer.
The union thus uses each ordinary layer either through two strands or not at all.
Across consecutive occupied ordinary layers, its edges form a perfect matching.
The two-choice cuts remain nonconsecutive, as in the thin-layer lemma.
Their number is at most $(D+1)/2$.
The extreme layers, marked incidences, and matching choices determine the union.
Consequently the number of these unions is at most
\[
 \bigl((D+1)^4+(D+1)^2\bigr)
 2^{3(6\Delta+4H+10)+(D+1)/2}.
\]
Since $\Delta,H\le C_\varepsilon$, this is at most
$B'_\varepsilon(b+1)^4 2^{b/2}$ for a constant $B'_\varepsilon$.

For a general normalized two-pole, use the series decomposition from
Theorem~\ref{thm:high-path-cycle-tradeoff}.
Every positive-rank piece has normalized path density at least
$1/2+\varepsilon$.
Let their ranks be $b_1,\ldots,b_q$, with sum $b$.
Packings contained in one piece contribute at most
\[
 B'_\varepsilon(b+1)^4\sum_i2^{b_i/2}
 \le2B'_\varepsilon(b+1)^4 2^{b/2}.
\]
Packings with one cycle in each of two pieces contribute at most
\[
 K_\varepsilon^2\sum_{i<j}2^{(b_i+b_j)/2}
 \le K_\varepsilon^2b^2 2^{b/2}.
\]
Increase the constant to cover both terms.
\end{proof}

\begin{corollary}[A cross-sum forest bound]
\label{cor:four-pole-cross-sum-forest}
For each $\eta>0$, there is a constant $A_\eta$ with the following property.
Let $K$ be a connected four-pole of rank $h$ with left ports $L_1,L_2$
and right ports $R_1,R_2$.
If
\[
 \sum_{i,j=1}^2P_{L_iR_j}(K)\ge(2+\eta)2^h,
\]
then the total number of four-terminal forests in $K$ is at most
\[
 A_\eta(h+3)^4 2^{h/2}.
\]
\end{corollary}

\begin{proof}
Add a new vertex $z$ adjacent to both left ports.
Add a new vertex $w$ adjacent to both right ports.
The resulting graph $N$ is a normalized two-pole of rank $h+2$.
Its terminal-path count is the displayed cross sum, so its normalized
density is at least $1/2+\eta/4$.

Append the four new edges to any four-terminal forest of $K$.
A through pair becomes one cycle containing $z,w$.
Two same-side caps become two vertex-disjoint cycles, one through each new vertex.
This operation is injective.
Apply Lemma~\ref{lem:high-path-two-cycle-packings} and absorb the factor two
from $2^{(h+2)/2}$ into $A_\eta$.
\end{proof}

\paragraph{Status.}
The exact state identity and the complete forest estimate passed the
path-stability agent's independent audit.
These estimates supply the state-counting input for
Theorem~\ref{thm:all-four-cut-half-exponent}.
That later theorem combines them with the aggregate inequality and the
separate analysis of bounded-rank shores.
